\section{задание: Переформатирование и извлечение сущностей}
\subsection{исходные данные}
\indent Входной текст (Тип документа): \textbf{Карточка товара маркетплейса} \\
\indent Поля для извлечения (JSON Schema): \textbf{\{product,price,currency,specs\{\}\}} \\
\indent Файл <<1.html>> содержит полную карточку товара. \\
\indent Файл <<2.html>> содержит повреждённую карточку товара. 
\subsection{базовый промпт}
\begin{verbatim}
					Верни JSON с ключами: product (строка), price (число),
					currency (строка), specs (словарь: строка:строка).
\end{verbatim}
Базовый промпт составил по примеру: ключи и их значения. \\
\noindent Системный промпт: 
\begin{verbatim}
					Ты - парсер данных. Твоя задача - только извлекать информацию.
					Не добавляй пояснений.
\end{verbatim}
Системный промпт из примера.
\subsection{журнал промпт-инжиниринга}
\noindent Промпт v1 + Системный промпт:
\begin{verbatim}
					Системный промпт:
						Ты - парсер данных.
						Твоя задача - только извлекать информацию.
						Не добавляй пояснений.
					Обычный промпт:
						Верни JSON с ключами: product (строка), price (min, число),
						currency (строка), specs (словарь: строка:строка).
\end{verbatim}
\newpage
\noindent Результат (YandexGPT, файл <<1.html>>):
\begin{figure}[H]
	\centering
	\includegraphics[width=15cm]{z6/a1.png}
\end{figure}
\noindent Анализ ошибок: модель верно извлекла JSON со странницы. \\
Корректировка: не требуется. \\
\newpage
\noindent Результат (DeepSeek, файл <<1.html>>):
\begin{figure}[H]
	\centering
	\includegraphics[width=15cm]{z6/b1.png}
\end{figure}
\noindent Анализ ошибок: модель вернула большую цену с карточки товара, потому что меньшая цена помечена спец. условиями. \\
Корректировка: не требуется. \\
\newpage
\noindent Результат (YandexGPT, файл <<2.html>>):
\begin{figure}[H]
	\centering
	\includegraphics[width=15cm]{z6/c1.png}
\end{figure}
\noindent Анализ ошибок: модель верно извлекла JSON со странницы, заменяя отсутствующие параметры на null. \\
Корректировка: не требуется. \\\newpage
\noindent Результат (DeepSeek, файл <<2.html>>):
\begin{figure}[H]
	\centering
	\includegraphics[width=15cm]{z6/d1.png}
\end{figure}
\noindent Анализ ошибок: модель верно извлекла JSON со странницы и достала где-то из глубин кода странницы валидную цену товара... \\
Корректировка: Лучше проверять исходные данные перед применением. \\
\subsection{Сравнение моделей:}
\begin{table}[h!]
	\centering
	\begin{tabular}{|l|p{5cm}|p{5cm}|}
		\hline
		\textbf{Критерий} & \textbf{YandexGPT} & \textbf{DeepSeek} \\
		\hline
		Синтаксическая валидность & 5 & 5 \\
		\hline
		Схемное  соответствие & 5 & 5 \\
		\hline
		\textbf{Итог} & 5 & 5 \\
		\hline
	\end{tabular}
\end{table}
Модель YandexGPT правильно отработала это задание, потому что искала значения в основном теле карточки товара (которая отображается пользователю). Но и модель DeepSeek технически верно выполнила задание, так как нашла верную цену товара в другой части кода карточки товара.
\newpage
\subsection{Финальный промпт:}
\begin{verbatim}
					Системный промпт:
						Ты - парсер данных.
						Твоя задача - только извлекать информацию.
						Не добавляй пояснений.
					Обычный промпт:
						Верни JSON с ключами: {ОПИСАНИЕ_JSON_СТРУКТУРЫ}.
\end{verbatim}
\subsection{Вывод по заданию:}
\indent В данном задании хорошо работают техники <<Ролевое моделирование>> и <<Обучение  на  примерах>>. Вместе они задают роль обработчика и обучают модель правильно выводить данные.